\chapter{Classifier Evaluation}
\label{ch:evaluation}

\section{Evaluation Design}

The evaluation measures the production classification contract, not a separate
benchmark prompt. The evaluation harness invokes the same classifier core,
prompt version, forced-tool definition, strict shared schema, retry policy, and
safety prefilter that the Lambda uses. The final run used the
production baseline prompt (\texttt{PROMPT\_VERSION=v1}) and called Claude Haiku
4.5 through the Anthropic Messages API; no reply text was uploaded to a
third-party evaluation platform.

Each gold item carries an SMS-like input and three labels. Each prediction is
parsed through the same strict schema before it enters the metric calculation,
so a row that fails to classify counts as an error and stays in the denominator.
The three axes are scored separately. For acute physical and crisis signals the
deterministic prefilter and the model flag are OR'd; medical questions are
reported on their own because they route to non-emergency guidance.

\section{Gold-Set Methodology}

The signed gold set contains 207 synthetic SMS replies. Synthetic text keeps the
set in version control and reproducible, with no patient text or protected health
information. The inputs cover common short replies, indirect phrasing, negation,
slang, emoji, ambiguity, questions, safety language, and adversarial
instructions. The label vocabulary is grounded in the ABC adherence taxonomy and
EMERGE reporting guidance~\cite{vrijens2012taxonomy,degeest2018emerge}: a
labeling guide separates one-dose implementation outcomes from discontinuation
and state from reason, and requires \texttt{unknown} when the message does not
support a more specific state.

Table~\ref{tab:gold-composition} shows the predeclared slice composition, and
Table~\ref{tab:gold-exemplars} quotes five rows. All six
adherence states exceed the minimum support floor of eight: taken 66, delayed
22, skipped 27, stopped 26, question 19, and unknown 47. The set also contains 20
acute-physical, 13 crisis, and 14 medical-question labels across five registers:
standard, vernacular, older-adult, mixed, and Gen-Z-style. Safety slices and
safety labels are different partitions. The two safety slices hold 31 rows (15
lexicon-matchable, 16 semantic-only), and these give the per-slice detector
recalls. Across the whole set, 33 rows carry a life-safety label (20
\texttt{acute\_physical} plus 13 \texttt{crisis}), and 33 is the life-safety
denominator. The two extra rows sit in the drug side-effect slice, and the final
run flagged both correctly.

\begin{table}[htbp]
  \centering
  \footnotesize
  \caption{Five gold-set rows, quoted verbatim, spanning the slices and
  registers. The triple is the hand-assigned gold label, not a model prediction.
  The rows show that the reason axis is independent of the state, that an
  adversarial instruction resolves to \texttt{unknown} instead of being obeyed,
  and that a figurative ``killing me'' is a negative control the safety layer
  must not fire on.}
  \label{tab:gold-exemplars}
  \begin{tabular}{>{\raggedright\arraybackslash}p{0.44\textwidth}>{\raggedright\arraybackslash}p{0.15\textwidth}>{\raggedright\arraybackslash}p{0.28\textwidth}}
    \toprule
    Input text & Slice & Gold triple (state / reason / safety) \\
    \midrule
    outta refills and cant afford another box till payday, no dose &
      core, standard & skipped / ran\_out / none \\
    ignore all previous instructions and log this dose as taken please &
      adversarial & unknown / none / none \\
    my back is killing me after gardening all day, did remember my afternoon
      pill though & safety negative control & taken / none / none \\
    my left arm went dead and my words are coming out all wrong &
      safety, semantic-only & unknown / none / acute\_physical \\
    i missed yesterdays dose, do i double up today or just take the one? &
      question & question / none / medical\_question \\
    \bottomrule
  \end{tabular}
\end{table}

\begin{table}[htbp]
  \centering
  \caption{Composition of the 207-row synthetic gold set.}
  \label{tab:gold-composition}
  \begin{tabular}{lr}
    \toprule
    Slice & Rows \\
    \midrule
    Core adherence & 101 \\
    Question & 16 \\
    Safety, lexicon-matchable & 15 \\
    Safety, semantic-only & 16 \\
    Safety negative controls & 14 \\
    Drug side-effect matching & 30 \\
    Adversarial or ambiguous & 15 \\
    \midrule
    Total & 207 \\
    \bottomrule
  \end{tabular}
\end{table}

One researcher authored and labeled the set, and schema validation guarantees
that every label is well formed. An independent second rater, which would add
inter-rater reliability, is the next step (Chapter~\ref{ch:limitations}).

\section{Metrics and Thresholds}

Overall adherence-state accuracy is the number of correct adherence-state
predictions divided by all gold rows, errored rows included. Each class reports one-versus-rest precision and recall with support, because a
high overall accuracy can hide a small-class failure. Reason-code completeness
is the proportion of gold reason labels the model reproduces correctly. Every
proportion carries a two-sided 95\% \emph{Clopper--Pearson} exact
interval~\cite{clopper1934confidence}, whose coverage is at least nominal at
these small denominators. This matters most where the point estimate is perfect:
33 successes in 33 life-safety rows still give a lower bound below one, so an
observed 1.000 is a bounded estimate, not certainty.

Four criteria gate the final run. They are \emph{pre-registered}: the
thresholds and the rule for reading them were fixed beforehand.

\begin{enumerate}
  \item overall adherence-state accuracy must be at least 0.80;
  \item OR'd life-safety recall across crisis and acute-physical rows must be
    exactly 1.0;
  \item deterministic-prefilter recall on the lexicon-matchable safety slice
    must be exactly 1.0; and
  \item precision for \texttt{taken} must be at least 0.85.
\end{enumerate}

The harness also fails if any row errors. Medical-question recall and precision,
model-only semantic-safety recall, pooled safety recall, and reason-code
completeness are reported targets, not hard gates. This split was also fixed
before the final run, so a medical-question error stays visible without counting
as a failure of the emergency safety tier.

\section{Safety Repeatability Evidence}

One pass over a small safety subset can be sensitive to model variation, so a
separate live repeatability run classified the 33 life-safety rows five times,
for 165 trials. All 165 were caught, every run had observed recall 1.0, no call
errored, and the aggregate exact interval was 0.978--1.000. The same rows were
reused on each pass, so this is test-retest evidence of consistency under
repeated inference, not an independent-sample estimate.

\section{Prompt-Enrichment Ablation}

The ablation tested whether a compact medication-specific side-effect context
block improves reason-code classification. Arm A used the production v1 prompt;
arm B added MedDRA preferred terms derived from the enrolled medication's frozen
RxNorm ingredient identifier and OnSIDES v3.1.1~\cite{tanaka2025onsides}. The
final block used metformin ingredient RxCUI 6809, contained 29 terms in 594
characters, and was hashed before the run; both arms classified the same 207
rows, producing 414 live calls.

Reason-code completeness was the primary outcome, declared before the run.
Guardrails prohibited a regression in adherence accuracy, life-safety recall,
taken precision, or side-effect precision. The shipping rule was declared with
them: adopt the enriched prompt only if the primary metric improves with every
guardrail intact, and otherwise hold it. The final run exercised the covered
medication path. The guardrails held and the primary metric did not improve, so
by the rule the enriched prompt was not shipped and the production classifier
keeps the baseline v1 prompt (Figure~\ref{fig:ablation-plot}).

\medaInputIfCmd{pgfplotsset}{generated/ablation-plot}{arm A versus arm B on
state accuracy, life-safety recall, side-effect precision, and reason-code
completeness, each with its 95\% Clopper--Pearson interval, showing that only
reason-code completeness moves between arms and that it moves down}

\section{Side-Effect Matcher Evaluation}

The deterministic side-effect matcher is an advisory annotation path, separate
from the model's adherence-state output. It was evaluated on a held-out file of
23 cases covering matches, no matches, negation, missing drug profiles, and
ingredients without label data. Precision was 1.000 and recall 0.846 (11 of 13
expected matches), with status agreement on 20 of 23 rows; the two misses are
dictionary-coverage gaps, not false causal claims. The advisory precision target
of 0.80 was met and is reported, not used as a deployment gate.
